\section{Canonical equality vacuity for the exact K230 core}
\label{sec:canonical-ve-transport}

The ordinary equality interpretation from Source~116 also applies to the
present K230 source.  It need not be routed through that source's older
CP/GIC adapter.  The essential observation is that the current ordinary
proof comparison in \texttt{CompleteBase230.v} and \texttt{Complete230.v}
can be parameterized by an interpreted theorem predicate.  The following
argument supplies the missing source-specific premise, while preserving
the coupling between the emitted matrix code and its private output label.
This is an ordinary mathematical proof, formalizable in ZFC by induction
on finite formulas and finite derivations.  It is not a claim that a new
exact-machine Rocq endpoint has been checked.

\subsection{The two numeric calculi}

Write $P(a,b)=\operatorname{tri}(a+b)+b$, and use the retained constructors
\[
 I(A,B)=4P(A,B),\quad M(a,b)=4P(a,b)+1,\quad
 Q_a A=4P(a,A)+2,\quad \mathsf{Eq}(a,b)=4P(a,b)+3.
\]
Here $F=1$, $N A=I(A,F)$, $A\mathbin{\&}B=N I(A,N B)$, and
$\exists_a A=N Q_a N A$.  Let $K(a,b,k,A)$ be the exact current core,
with input label $a$, output label $b$, private container label $k$,
and matrix code $A$.  In particular,
\begin{align*}
 S(a,b,k)&=Q_k I(M(k,a),M(k,b)),\\
 R(b,k,A)&=Q_b I\bigl(A,\exists_b(A\mathbin{\&}M(b,k))\bigr),\\
 K(a,b,k,A)&=\exists_k\Bigl(
 I(A,M(b,k))\mathbin{\&}
 Q_b I\bigl(M(b,k),\exists_a(A\mathbin{\&}
 Q_a I(S(a,b,k),R(b,k,A)))\bigr)\Bigr).
\end{align*}
These expressions retain the repeated and shadowed binders of the source.

The old enabled bundle at $(a,b,c)$ is
\[
 K(a,b,c,c)\mathbin{\&}
 \bigl(I(M(a,b),Q_cM(a,b))\mathbin{\&}
 I(\mathsf{Eq}(a,b),Q_c\mathsf{Eq}(a,b))\bigr).
\]
The proposed replacement changes only the final component to
\[
 I\bigl(\mathsf{Eq}(a,b),Q_{\mathsf{Eq}(a,b)}\mathsf{Eq}(a,b)\bigr).
 \tag{CVE}
\]
The gate remains $c\ne a,b$.  Both calculi retain numeric truth $0$,
all eight original logical rows, modus ponens, and closed-theorem
generalization.  Denote them by $H_{\rm old}$ and $H_{\rm cve}$.
Classically equivalent DN4 parenthesizations of the displayed bundles
give the same argument, by quantified congruence.

\paragraph{The easy direction.}
For every $a,b$, the label
$k=\mathsf{Eq}(a,b)=4P(a,b)+3$ is larger than both endpoints.  Projecting
the old bundle at $(a,b,k)$ supplies (CVE).  The old bundle at $(a,b,c)$
supplies its unchanged core and membership-vacuity components.  Classical
conjunction introduction reconstructs every new enabled bundle.  Thus
every $H_{\rm cve}$ derivation is an $H_{\rm old}$ derivation, and
\[
 \operatorname{Con}(H_{\rm old})
 \ \Longrightarrow\ \operatorname{Con}(H_{\rm cve}).
\]
Conversely, the new calculus really supplies (CVE) for every $a,b$:
choose any $c\ne a,b$ and project the enabled bundle.  Its membership
vacuity is still available at every gated triple.  Retaining the original
eighth logical row supplies classicality independently of these projections.

\subsection{The equality-profile interpretation}

For $z\ne a,b$, let
\begin{align*}
 L_{ab}(z)&=(M(z,a)\to M(z,b))\mathbin{\&}
             (M(z,b)\to M(z,a)),\\
 U_{ab}(z)&=(M(a,z)\to M(b,z))\mathbin{\&}
             (M(b,z)\to M(a,z)),\\
 E_{ab}^{z}&=(\forall_z L_{ab}(z))\mathbin{\&}
             (\forall_z U_{ab}(z)),\\
 E(a,b)&=E_{ab}^{\mathsf{Eq}(a,b)}.
\end{align*}
These are finite named formulas; the arrows, conjunctions and negations
are subsequently encoded with the retained numeric constructors.
Let $J$ preserve membership, implication, primitive negation and every
universal binder, replacing only the equality atom $(a=b)$ by $E(a,b)$.

The uniform bootstrap in Source~116, Sections~1--3 of
\texttt{116-eqred-audit-mathematical\_audit.md}, applies to these exact
retained logical rows.  It derives all interpreted logical axioms from
full membership vacuity and (CVE).  In particular it handles equality
witnesses, membership congruence, transitivity, and full vacuity for the
\emph{interpreted} equality $E$.  Auxiliary labels are transported through
a fresh maximum-plus-one label.  This includes coincidences with an
endpoint's canonical auxiliary label and includes label~$0$.

Consequently every theorem $p$ of the full pure named logic has a finite
$H_{\rm cve}$ proof of $\operatorname{enc}(Jp)$.  This theorem is obtained
by induction over the logical axiom, modus-ponens and closed-GEN
constructors.  It does not assert that primitive full equality vacuity
is derivable, and it does not assume that $J$ commutes with substitution
or with a syntactic renaming operation.

Define the theorem predicate
\[
 \mathcal Q_E(p)\quad\Longleftrightarrow\quad
 H_{\rm cve}\vdash\operatorname{enc}(Jp).
\]
It is closed under pure logical consequences, modus ponens, and
generalization.  These are exactly the properties used by the
\texttt{Q}-based infrastructure of \texttt{CompleteBase230.v}.

\subsection{Supplying the ordinary K core with an actual source code}

Fix an arbitrary ordinary matrix $\varphi$.  Use precisely the current
translation $p=\texttt{tr source\_names}\,\varphi$, with named output~$3$,
input~$1$, and the retained ordinary treatment of bottom.  Let
\[
 c=\operatorname{enc}(\neg\neg p),\qquad
 c'=\operatorname{enc}(\neg\neg Jp).
\]
The original code $c$ is exactly
\texttt{ordinary\_code230}\,$\varphi$.  Both $c,c'>3$.  More strongly,
both exceed every label occurring anywhere in $p$, including bound labels.
For $c$ this is the finite-tree coding bound.  For $c'$ use the same
bound and the fact that every original label is retained somewhere in
$Jp$.  Each equality expansion uses both endpoints and only adds its
auxiliary label.  No decoder inverse is required for this argument.

Request the \emph{actual} enabled new bundle at $(1,3,c')$.  Its matrix
argument is the number $c'$ itself, so the constructor equations give
literally
\[
 K(1,3,c',c')
  =\operatorname{enc}\bigl(\mathsf K(1,3,c',\neg\neg Jp)\bigr),
\]
where $\mathsf K$ is the displayed named version of the same core.
Projection and double-negation equivalence, propagated under every
connective and quantifier, therefore give
\[
 H_{\rm cve}\vdash
 \operatorname{enc}\bigl(\mathsf K(1,3,c',Jp)\bigr).
 \tag{1}
\]
Since the core contains no equality atom outside its matrix,
\[
 J\bigl(\mathsf K(1,3,c',p)\bigr)
   =\mathsf K(1,3,c',Jp).
 \tag{2}
\]

In the full ordinary named logic,
\[
 \mathsf K(1,3,c',p)\ \longleftrightarrow\
 \mathsf K(1,3,c,p)
 \tag{3}
\]
is a pure alpha-equivalence theorem.  Both private labels are absent
from $p$ and differ from $1,3$.  Rename the outer container binder and
the locally shadowing subset binder; the matrix and the input/output
binders remain literal.  Equivalently their de Bruijn translations
under an injective environment are equal, so the retained
\texttt{Std\_alpha.alpha\_dB} theorem proves (3).
If $c=c'$, (3) is identity.

Apply the already established pure logical interpretation to (3), and
use (1)--(2).  The result is
\[
 \boxed{\mathcal Q_E\bigl(
   \mathsf K(1,3,\texttt{ordinary\_code230}\,\varphi,
                 \texttt{tr source\_names}\,\varphi)\bigr).}
 \tag{4}
\]
This is the exact counterpart of
\texttt{Complete230.ordinary\_named\_core\_Q}.
The reasoning never obtains a source instance with an independently
chosen output coordinate.  It obtains the genuine $(c',c')$ instance
first and derives the desired named theorem afterward.
Possible collisions between $c$ and an auxiliary binder introduced
inside $Jp$ cause no gap: the alpha theorem is proved for $p$ first and
then interpreted as a whole.  No commuting law between $J$ and renaming
is used.

\subsection{Transport of the ordinary inconsistency proof}

Replace the original predicate $Q$ by $\mathcal Q_E$ in the mathematical
argument of \texttt{CompleteBase230.v} and \texttt{Complete230.v}.
Define $TB_E$ to consist of closed ordinary readings
$\operatorname{Lib}(n,p)$ whose named formula satisfies
$\mathcal Q_E(p)$, with the same bound on free labels as before.
Then:
\begin{enumerate}
\item The equality axioms belong to $TB_E$ because they are pure logical
      theorems before interpretation.
\item Ordinary ND proof discharge and the closed-formula round trip
      use only interpreted pure logical theorems, MP and GEN.
      In particular the original use of \texttt{alpha\_dB} is interpreted;
      no new syntactic commutation premise is introduced.
\item Equation~(4), the unchanged
      \texttt{named\_source\_to\_kprime} argument, and finite proof
      collection show that every ordinary K-prime axiom is provable from
      $TB_E$.
\item An ordinary contradiction yields
      $\mathcal Q_E(\neg\forall_0(0=0))$.
      Interpreted pure reflexivity and GEN also yield
      $\mathcal Q_E(\forall_0(0=0))$.  Their numeric modus ponens gives
      exactly $F=1$.
\end{enumerate}
The optional irreflexivity premise used by the existing bridge is still
available: its named representative is $\neg M(0,0)$, unchanged by $J$,
and its code is literally $I(F,F)$.

The retained K-prime-to-ZFC finite-proof comparison now gives
\[
 \mathrm{ZFC}\vdash\bot\quad\Longrightarrow\quad
 H_{\rm cve}\vdash F.
\]
Combining this implication with the easy direction and the existing
K230 consistency calibration proves, in the ordinary syntactic sense,
\[
 \operatorname{Con}(H_{\rm cve})
   \quad\longleftrightarrow\quad\operatorname{Con}(\mathrm{ZFC}).
\]
The bootstrap derivations, fresh-label choices, formula translations and
finite-proof maps are all explicit operations on finite objects; their
correctness inductions can be carried out in ZFC.  A machine-state claim
still requires a correctly compiled literal source with the indicated
row coverage, a full runtime certificate, and the blank-input simulation
proof.  This section alone supplies no smaller transition table.

\paragraph{Code-zero and relative-falsehood boundary.}
The argument uses constructor homomorphism on finite encoded trees,
not the false statement $\operatorname{enc}(D(n))=n$.  The source is
requested at the explicitly constructed positive code $c'$, never at
a recoding of an arbitrary decoder output.  The numeric falsehood is
always $1$; under a logical negation it is fixed as the designated
bottom.  The equality-profile bootstrap treats membership occurrences
that collapse to $M(0,0)$ using the retained congruence and explosion
derivations.  Thus neither the native truth-code alias at $0$ nor the
relative membership interpretation of $1$ is silently replaced by an
unrestricted raw-syntax renaming principle.
